% PLP Selection Examination — Part B (practical)
% Layout follows the ISI Computing Laboratory lab-test papers.
% Build:  pdflatex docs/exam-paper.tex
% Keep in step with scenarios/*/ticket.md — the VM serves those, this is print.
\documentclass[11pt,a4paper]{article}
\usepackage[margin=1in,top=0.9in,bottom=1in]{geometry}
\usepackage{enumitem}
\usepackage{framed}
\usepackage{xcolor}
\usepackage{fancyvrb}
\usepackage{array}
\usepackage[T1]{fontenc}

\pagestyle{plain}
\setlength{\parindent}{0pt}
\setlength{\parskip}{0.45em}

% shell listings, as the sample paper sets sample input/output
\DefineVerbatimEnvironment{shell}{Verbatim}{fontsize=\small,xleftmargin=1.5em}

\newcommand{\hd}[1]{\medskip\textbf{#1}\par\nopagebreak}

\begin{document}

%----------------------------------------------------------------- header
\begin{center}
  {\large\textbf{INDIAN STATISTICAL INSTITUTE}}\\[0.6em]
  Project Linked Person --- Selection \qquad 2026--2027\\[0.35em]
  Subject: Systems Support (Software, Hardware and Web)\\[0.35em]
  Part B --- Practical Examination\\[0.35em]
  Total: $16 + 18 + 16 = 50$ marks \qquad Duration: 90 minutes
\end{center}

\begin{framed}
\begin{center}\textbf{EXAMINATION INSTRUCTIONS}\end{center}
\begin{enumerate}[leftmargin=1.5em,itemsep=0.35em,topsep=0.3em]
  \item You have root on a virtual machine. It is a disposable clone: you
        cannot damage anything that matters, and it is connected to no real
        network. Log in with
        \verb|ssh candidate@192.168.56.10|.
  \item Each question has four parts --- \emph{the problem}, \emph{how to
        reproduce it}, \emph{what to fix}, and \emph{how to check your work}.
        Run the checks yourself; you should not have to guess whether you are
        finished.
  \item \textcolor{red}{\textbf{The \emph{Constraints} of every question are
        marked. Breaking a constraint loses marks even when the symptom goes
        away.} Examples: \texttt{chmod 777}, pinning names into
        \texttt{/etc/hosts}, reformatting a filesystem.}
  \item Keep \verb|~/FIXLOG.md| up to date as you go: symptom, cause, what you
        changed, and how you verified it. It is marked, and there is a short
        oral discussion of your work afterwards.
  \item You are not expected to know all of this in advance. Partial credit is
        real: a fault you diagnose correctly and explain, but cannot finish,
        is worth more than one you paper over.
\end{enumerate}
\end{framed}

\textbf{NOTE:} \texttt{man}, \verb|--help|, \texttt{journalctl} and the logs
under \texttt{/var/log} are all available on the machine. Your commands are
recorded. Questions may be attempted in any order.

%----------------------------------------------------------------- questions
\begin{enumerate}[label=\textbf{Q\arabic*.},leftmargin=2.6em,labelsep=0.6em,itemsep=1.6em,topsep=1.2em]

%================================================================== Q1
\item \textbf{Disk full on the lab file server.} \hfill [16 marks]

\emph{Reported by a research scholar:} ``The reporting log stopped updating
before 07:00. I cannot copy anything into \texttt{/srv/data} --- it says the
disk is full. Submitting a job into \texttt{/srv/spool} \emph{also} says the
disk is full, but \texttt{df -h} shows that one has plenty of space left. That
part makes no sense to me.''

Two storage faults, on two different filesystems. \texttt{/srv/data} is
genuinely full, so the reporting daemon can no longer write. \texttt{/srv/spool}
refuses to create files \emph{although \texttt{df} reports free space} ---
something other than free bytes has run out.

\hd{How to reproduce it}
\begin{shell}
df -h /srv/data
tail -1 /srv/data/logs/reportd.log     # last entry is hours old
df -h /srv/spool                       # shows space free
touch /srv/spool/incoming/testjob      # "No space left on device"
sudo du -sh /srv/data                  # smaller than df says is in use
\end{shell}

\hd{What to fix}
\begin{enumerate}[label=\textbf{\arabic{enumi}.\arabic*},leftmargin=2.8em,itemsep=0.3em]
  \item Bring \texttt{/srv/data} back to \textbf{under 10\% used} (the
        \texttt{Use\%} column of \texttt{df -h}). Freeing the largest obvious
        files will not get you there on its own. Treat any difference between
        what \texttt{df} says is used and what \texttt{du} can account for as
        something you still have to explain.
  \item Get the reporting daemon writing to
        \texttt{/srv/data/logs/reportd.log} again. It is running; its writes
        were failing. Check the last line's timestamp rather than assuming.
  \item Make \texttt{/srv/spool} able to accept new files. It is a separate
        filesystem with a separate fault and it is \emph{not} short of space:
        freeing bytes there will not help. A filesystem can exhaust more than
        one kind of resource, and \texttt{df} has another mode that reports the
        other one.
  \item In \verb|~/FIXLOG.md|, explain in two or three sentences why
        \texttt{/srv/spool} reported ``No space left on device'' while
        \texttt{df} showed free space.
\end{enumerate}

\hd{How to check your work}
\begin{shell}
df -h /srv/data                       # Use% under 10%
sudo systemctl restart reportd; sleep 5
tail -1 /srv/data/logs/reportd.log    # timestamp within the last few seconds
touch /srv/spool/incoming/t && echo OK && rm /srv/spool/incoming/t
ls /srv/data/current                  # five survey-block CSV files, still there
\end{shell}

\hd{Constraints --- these are marked}
\begin{itemize}[leftmargin=1.4em,itemsep=0.2em]
  \item Nothing under \texttt{/srv/data/current} may be deleted or changed.
        It is live survey data with no second copy on this machine.
  \item Do not reformat, recreate, resize or unmount either filesystem.
  \item \texttt{/srv/data/archive/README.txt} says what is safe to remove
        there. Read it before deleting anything.
  \item Rebooting clears one of these faults. On a real file server at 11am
        you would not have that option; treat it as unavailable.
\end{itemize}

%================================================================== Q2
\item \textbf{Shared project folder unusable.} \hfill [18 marks]

\emph{Reported for the Statistics Lab group:} ``Bikram cannot save anything
into \texttt{shared/}, and cannot open \texttt{notes.md} at all. When I create
a file there, the others cannot edit it --- it comes out belonging to my own
personal group instead of the project group.
\texttt{sudo systemctl restart reportd} used to work for all of us; now it is
refused. And Chandan from the internal audit cell needs to \emph{read} that
folder for the annual review. He cannot see anything.''

\texttt{/srv/projects/statlab/} is a shared area for the \texttt{statlab}
group. Several things about it have been changed, and the auditor has no
access at all.

\hd{How to reproduce it}
\begin{shell}
sudo -u bikram touch /srv/projects/statlab/shared/test     # Permission denied
sudo -u bikram cat /srv/projects/statlab/shared/notes.md   # Permission denied
sudo -u anita sudo -n systemctl restart reportd            # refused
sudo -u chandan cat /srv/projects/statlab/shared/notes.md  # Permission denied

stat -c '%a %U:%G %n' /srv/projects/statlab
id bikram ; getfacl /srv/projects/statlab
sudo cat /etc/sudoers.d/statlab
\end{shell}

\hd{What to fix}
\begin{enumerate}[label=\textbf{\arabic{enumi}.\arabic*},leftmargin=2.8em,itemsep=0.25em]
  \item Everyone in \texttt{statlab} can read and write everything under
        \texttt{/srv/projects/statlab/}.
  \item A file created there by one member ends up usable by the whole group.
  \item \texttt{notes.md} is a shared file again, not one person's private file.
  \item Members of \texttt{statlab} --- and only them --- can run
        \texttt{sudo -n systemctl restart reportd} without a password.
  \item Chandan can \textbf{read} that tree.
\end{enumerate}

\hd{How to check your work}
Test as the affected users, not as root.
\begin{shell}
sudo -u bikram touch /srv/projects/statlab/shared/t1 && echo "write OK"
stat -c '%U:%G' /srv/projects/statlab/shared/t1    # group must be: statlab
sudo rm -f /srv/projects/statlab/shared/t1
sudo -u bikram cat /srv/projects/statlab/shared/notes.md >/dev/null && echo "read OK"
sudo -u anita sudo -n systemctl restart reportd && echo "sudo OK"
sudo -u chandan cat /srv/projects/statlab/shared/notes.md >/dev/null && echo "auditor read OK"
sudo -u chandan touch /srv/projects/statlab/shared/t2 && echo "PROBLEM: auditor can write"
\end{shell}
The second line is the one people miss. If the group comes out as anything
other than \texttt{statlab}, the second complaint is not fixed yet.

\hd{Constraints --- these are marked}
\begin{itemize}[leftmargin=1.4em,itemsep=0.2em]
  \item \textbf{Chandan must not be put into the \texttt{statlab} group.}
        Audit policy forbids putting auditors into the groups they audit. He
        must be able to read and must \emph{not} be able to write.
  \item No \texttt{chmod 777}, anywhere. This is a multi-user machine.
  \item Do not change anyone's login shell, password or home directory.
\end{itemize}

\emph{Hint.} One of these faults is about which \textbf{bits are set on a
directory}, not about who owns it --- a directory carries more than the nine
permission bits you see first. Group membership changes do not affect a shell
that is already open. Edit anything under \texttt{/etc/sudoers.d/} with
\texttt{visudo -f}.

%================================================================== Q3
\item \textbf{Campus names stopped resolving.} \hfill [16 marks]

\emph{Reported by the systems group:} ``This host cannot resolve any campus
name since the patching window. \texttt{ping www.isi.local} sits there for
several seconds and then fails. The odd one is \texttt{portal.isi.local} ---
that \emph{does} answer, but with an address that has not been ours since 2023.
Please do not paper over this with hosts-file entries; we have forty more
machines behind this one.''

The host runs its own resolver (\textbf{dnsmasq}) serving the
\texttt{isi.local} zone and forwarding everything else upstream. The addresses
it is supposed to hand out:

\begin{center}
\begin{tabular}{>{\ttfamily}l @{\qquad} >{\ttfamily}l}
\textrm{\textbf{name}} & \textrm{\textbf{address}}\\[0.15em]
www.isi.local     & 127.0.0.1\\
statlab.isi.local & 127.0.0.1\\
portal.isi.local  & 127.0.0.1\\
git.isi.local     & 10.10.10.9\\
nas.isi.local     & 10.10.10.20\\
\end{tabular}
\end{center}

\hd{How to reproduce it}
\begin{shell}
getent hosts www.isi.local        # nothing, after a long pause
getent hosts git.isi.local        # nothing
getent hosts portal.isi.local     # answers 192.0.2.77 -- wrong address

dig +short @127.0.0.1 git.isi.local
systemctl status dnsmasq
\end{shell}

\hd{What to fix}
\begin{enumerate}[label=\textbf{\arabic{enumi}.\arabic*},leftmargin=2.8em,itemsep=0.25em]
  \item All five names resolve to the correct addresses, for every user on the
        host, in under a second.
  \item The answers come from the resolver, not from a local file.
  \item It still works after a reboot.
\end{enumerate}

\hd{How to check your work}
\begin{shell}
for n in www statlab portal git nas; do
  printf '%-22s %s\n' "$n.isi.local" "$(getent hosts $n.isi.local | awk '{print $1}')"
done
dig +short @127.0.0.1 git.isi.local     # must print 10.10.10.9
grep -i isi.local /etc/hosts            # must print nothing at all
systemctl is-active dnsmasq ; systemctl is-enabled dnsmasq
\end{shell}

\hd{Constraints --- these are marked}
\begin{itemize}[leftmargin=1.4em,itemsep=0.2em]
  \item \textbf{Do not put \texttt{isi.local} names into \texttt{/etc/hosts}.}
        The systems group will not maintain hosts files across forty machines.
        It makes the symptom disappear and loses marks.
  \item Do not make \texttt{/etc/resolv.conf} immutable
        (\texttt{chattr +i}) to hold a fix in place.
\end{itemize}

\emph{Hint.} There is more than one fault and they are in different layers. A
useful order to think in: \emph{is a DNS query being sent at all? where is it
being sent? is anything answering? is the answer correct?} A service that
refuses to start will usually say why, in its own log. The long pause before
failure is itself a clue about which layer is wrong.

\end{enumerate}
\end{document}
